% ============================================================
\chapter{Radon-Nikodym Theorem}
\label{ch:radon_nikodym}
% ============================================================

\section{Introduction}

This chapter establishes one of the most important results in measure theory: the
Radon-Nikodym theorem\index{Radon-Nikodym theorem}, which characterizes measures
that are absolutely continuous with respect to a given measure. This theorem is
fundamental in probability (densities, conditional expectation) and functional analysis.

\begin{intuition}[title=\textbf{Radon-Nikodym: changing the reference measure}]
If $\nu$ is absolutely continuous with respect to $\mu$, it means that $\nu$ only
``sees'' what $\mu$ sees: sets negligible for $\mu$ are also negligible for $\nu$.
The Radon-Nikodym theorem then asserts that $\nu$ is obtained by ``weighting'' $\mu$
with a density: $d\nu = f\, d\mu$. This generalizes the fact that any probability
on $\R$ with a density can be written as $\PP(A) = \int_A f(x)\, dx$.
\end{intuition}

% ------------------------------------------------------------
\section{Absolute continuity}
% ------------------------------------------------------------

\begin{definition}[Absolute continuity]
\label{def:abs_cont}
\index{absolute continuity}
Let $\mu$ be a positive measure and $\nu$ a signed measure on $(X, \mathcal{A})$.
We say that $\nu$ is \emph{absolutely continuous} with respect to $\mu$, written
$\nu \ll \mu$\index{$\ll$}, if:
\[
    \forall A \in \mathcal{A}, \quad \mu(A) = 0 \implies \nu(A) = 0.
\]
\end{definition}

\begin{proposition}[$\varepsilon$-$\delta$ characterization]
\label{prop:eps_delta}
If $\mu$ is a positive measure and $\nu$ is a \emph{finite} signed measure, then
$\nu \ll \mu$ if and only if:
\[
    \forall \varepsilon > 0,\, \exists\, \delta > 0,\, \forall A \in \mathcal{A},\quad
    \mu(A) < \delta \implies \abs{\nu(A)} < \varepsilon.
\]
\end{proposition}

\begin{proof}
($\Leftarrow$) Clear: if $\mu(A) = 0$, then $\mu(A) < \delta$ for every $\delta$,
giving $\abs{\nu(A)} < \varepsilon$ for every $\varepsilon$, hence $\nu(A) = 0$.

($\Rightarrow$) By contradiction. Suppose $\nu \ll \mu$ but the $\varepsilon$-$\delta$
condition fails. There exist $\varepsilon_0 > 0$ and a sequence $(A_n)$ with
$\mu(A_n) < 2^{-n}$ and $\abs{\nu(A_n)} \geq \varepsilon_0$. Set
$B_n = \bigcup_{k \geq n} A_k$. Then $\mu(B_n) \leq 2^{1-n}$ and
$B = \bigcap_n B_n$ satisfies $\mu(B) = 0$. By decreasing continuity,
$\abs{\nu}(B) = \lim_n \abs{\nu}(B_n) \geq \varepsilon_0 > 0$, contradicting
$\nu \ll \mu$.
\end{proof}

\begin{example}
\begin{enumerate}
    \item If $f \geq 0$ is measurable and $\nu(A) = \int_A f\, d\mu$, then $\nu \ll \mu$.
    \item The Dirac measure $\delta_0$ is not absolutely continuous with respect to
    Lebesgue measure $\lambda$: $\lambda(\{0\}) = 0$ but $\delta_0(\{0\}) = 1$.
    \item If $\nu \ll \mu$ and $\mu \ll \nu$, we say $\mu$ and $\nu$ are
    \emph{equivalent}, written $\mu \sim \nu$.
\end{enumerate}
\end{example}

% ------------------------------------------------------------
\section{The Radon-Nikodym theorem}
% ------------------------------------------------------------

\begin{theorem}[Radon-Nikodym]
\label{thm:radon_nikodym}
\index{Radon-Nikodym theorem}
Let $(X, \mathcal{A}, \mu)$ be a $\sigma$-finite measure space and let $\nu$ be a
$\sigma$-finite signed measure on $(X, \mathcal{A})$ with $\nu \ll \mu$. Then there
exists a measurable function $f : X \to [-\infty, +\infty]$, integrable with respect
to $\mu$ (if $\nu$ is finite), such that:
\[
    \nu(A) = \int_A f\, d\mu, \qquad \forall A \in \mathcal{A}.
\]
Moreover, $f$ is unique $\mu$-almost everywhere. If $\nu$ is a positive measure, then
$f \geq 0$ $\mu$-a.e.
\end{theorem}

\begin{proof}[Proof (case $\mu, \nu$ finite positive measures)]
We present Von Neumann's proof. Set $\lambda = \mu + \nu$, a finite positive measure.
For every $g \in L^2(\lambda)$, the linear functional
$\varphi(g) = \int g\, d\nu$ is continuous on $L^2(\lambda)$, since by
Cauchy--Schwarz:
\[
    \abs{\varphi(g)} \leq \int \abs{g}\, d\lambda
    \leq \bigl(\textstyle\int g^2\, d\lambda\bigr)^{1/2} \lambda(X)^{1/2}.
\]
By the Riesz representation theorem for Hilbert spaces, there exists
$h \in L^2(\lambda)$ such that:
\[
    \int g\, d\nu = \int gh\, d\lambda, \qquad \forall g \in L^2(\lambda).
\]
Taking $g = \ind_A$ gives $\nu(A) = \int_A h\, d\lambda$ for all
$A \in \mathcal{A}$. Since $0 \leq \nu(A) \leq \lambda(A)$, we deduce
$0 \leq h \leq 1$ $\lambda$-a.e.

We can write $\nu(A) = \int_A h\, d\mu + \int_A h\, d\nu$, hence
$\int_A (1-h)\, d\nu = \int_A h\, d\mu$. Setting $E_0 = \{h = 1\}$, one checks
$\mu(E_0) = 0$. Define $f = h/(1-h)$ on $E_0^c$ and $f = 0$ on $E_0$. The monotone
convergence theorem applied to $g_n = (1 + h + \cdots + h^n)\ind_A$ yields:
\[
    \nu(A) = \int_A f\, d\mu.
\]
Uniqueness of $f$ $\mu$-a.e.\ follows from the fact that if $\int_A f\, d\mu =
\int_A g\, d\mu$ for all $A$, then $f = g$ $\mu$-a.e.
\end{proof}

\begin{definition}[Radon-Nikodym derivative]
\index{Radon-Nikodym derivative}
The function $f$ from Theorem~\ref{thm:radon_nikodym} is called the
\emph{Radon-Nikodym derivative} of $\nu$ with respect to $\mu$ and is denoted:
\[
    f = \frac{d\nu}{d\mu}.
\]
\end{definition}

\begin{proposition}[Properties of the Radon-Nikodym derivative]
\label{prop:rn_props}
Let $\mu, \nu, \lambda$ be $\sigma$-finite measures.
\begin{enumerate}
    \item \textbf{Chain rule:} if $\nu \ll \mu \ll \lambda$, then
    $\dfrac{d\nu}{d\lambda} = \dfrac{d\nu}{d\mu} \cdot \dfrac{d\mu}{d\lambda}$
    $\lambda$-a.e.
    \item \textbf{Inverse:} if $\nu \ll \mu$ and $\mu \ll \nu$, then
    $\dfrac{d\mu}{d\nu} = \left(\dfrac{d\nu}{d\mu}\right)^{-1}$ $\mu$-a.e.
    \item \textbf{Linearity:} if $\nu_1, \nu_2 \ll \mu$, then
    $\dfrac{d(\alpha\nu_1 + \beta\nu_2)}{d\mu} = \alpha\dfrac{d\nu_1}{d\mu}
    + \beta\dfrac{d\nu_2}{d\mu}$.
\end{enumerate}
\end{proposition}

\begin{keyformulas}[title=\textbf{Key Radon-Nikodym formulas}]
If $\nu \ll \mu$ with $f = \frac{d\nu}{d\mu}$, then for every measurable $g \geq 0$
or $g \in L^1(\nu)$:
\[
    \int g\, d\nu = \int g \cdot f\, d\mu = \int g \cdot \frac{d\nu}{d\mu}\, d\mu.
\]
Chain rule: $\dfrac{d\nu}{d\lambda} = \dfrac{d\nu}{d\mu} \cdot \dfrac{d\mu}{d\lambda}$.
\end{keyformulas}

% ------------------------------------------------------------
\section{Lebesgue decomposition theorem}
% ------------------------------------------------------------

\begin{theorem}[Lebesgue Decomposition]
\label{thm:lebesgue_decomp}
\index{Lebesgue decomposition}
Let $(X, \mathcal{A}, \mu)$ be a $\sigma$-finite measure space and let $\nu$ be a
$\sigma$-finite signed measure. Then there exists a unique pair of signed measures
$(\nu_a, \nu_s)$ such that:
\[
    \nu = \nu_a + \nu_s, \qquad \nu_a \ll \mu, \qquad \nu_s \perp \mu.
\]
Moreover, $\nu_a(A) = \int_A f\, d\mu$ for some measurable $f$ (the Radon-Nikodym
derivative of $\nu_a$ with respect to $\mu$).
\end{theorem}

\begin{proof}[Proof sketch]
It suffices to treat the case where $\nu$ is a finite positive measure and $\mu$ is
finite. For each $\lambda > 0$, the signed measure $\nu - \lambda\mu$ admits a Hahn
decomposition: $X = P_\lambda \cup N_\lambda$ with $P_\lambda$ positive and
$N_\lambda$ negative for $\nu - \lambda\mu$. One may choose $P_\lambda$ increasing in
$\lambda$ (replacing $P_\lambda$ by $\bigcup_{\lambda' \leq \lambda} P_{\lambda'}$ if
needed). Set $\lambda_0 = \sup\{\lambda > 0 : \nu(N_\lambda) > 0\}$ and take a
sequence $\lambda_n \nearrow \lambda_0$. The set $N = \bigcap_n N_{\lambda_n}$
satisfies $\mu(N) = 0$: indeed, for every $n$ and every $A \subseteq P_{\lambda_n}^c$,
$\nu(A) \leq \lambda_n \mu(A)$, and passing to the limit gives $\mu(N) = 0$.
Set $\nu_s(\cdot) = \nu(\cdot \cap N)$ and $\nu_a(\cdot) = \nu(\cdot \cap N^c)$.
Since $\mu(N) = 0$, we have $\nu_s \perp \mu$. To show $\nu_a \ll \mu$: if
$\mu(A) = 0$ with $A \subseteq N^c$, then $A \subseteq P_{\lambda_n}$ for $n$ large
enough, and $\nu(A) \leq \lambda_n \mu(A) = 0$. The Radon-Nikodym theorem applied to
$\nu_a$ provides the density $f$.
See \textsc{Rudin}, \emph{Real and Complex Analysis}, Theorem~6.10.
\end{proof}

% ------------------------------------------------------------
\section{Applications to probability}
% ------------------------------------------------------------

\begin{example}[Probability density]
Let $X$ be a real random variable on $(\Omega, \mathcal{F}, \PP)$ whose law $\PP_X$
is absolutely continuous with respect to Lebesgue measure $\lambda$. Then there exists
$f_X \geq 0$ with $\int f_X\, d\lambda = 1$ such that:
\[
    \PP(X \in A) = \int_A f_X(x)\, dx, \qquad \forall A \in \mathcal{B}(\R).
\]
The function $f_X = \frac{d\PP_X}{d\lambda}$ is the \emph{density} of $X$.
\end{example}

\begin{example}[Conditional expectation]
Let $(\Omega, \mathcal{F}, \PP)$ be a probability space,
$\mathcal{G} \subset \mathcal{F}$ a sub-$\sigma$-algebra, and $Y \in L^1(\PP)$.
The signed measure $\nu(A) = \int_A Y\, d\PP$ for $A \in \mathcal{G}$ is absolutely
continuous with respect to $\PP|_{\mathcal{G}}$. By Radon-Nikodym, there exists a
unique (a.s.) $\mathcal{G}$-measurable function $Z \in L^1(\PP)$ with:
\[
    \int_A Z\, d\PP = \int_A Y\, d\PP, \qquad \forall A \in \mathcal{G}.
\]
This function $Z$ is the \emph{conditional expectation} $\E[Y | \mathcal{G}]$.
\end{example}

\begin{example}[Likelihood ratio]
Let $\PP$ and $\Q$ be two probability measures on $(\Omega, \mathcal{F})$ with
$\Q \ll \PP$. The \emph{likelihood ratio} is $L = \frac{d\Q}{d\PP}$, and for every
$\Q$-integrable random variable $Y$:
\[
    \E_{\Q}[Y] = \E_{\PP}[Y \cdot L].
\]
This is the foundation of change of measure in finance (Girsanov's theorem) and
statistics (Neyman--Pearson test).
\end{example}

% ------------------------------------------------------------
\section{Exercises}
% ------------------------------------------------------------

\begin{exercise}
Let $\mu$ be Lebesgue measure on $[0,1]$ and $\nu(A) = \int_A 2x\, dx$. Verify
that $\nu \ll \mu$, identify $\frac{d\nu}{d\mu}$, and compute
$\int_{[0,1]} x^2\, d\nu$.
\end{exercise}

\begin{exercise}
Show that if $\nu \ll \mu$ and $\nu \perp \mu$, then $\nu = 0$.
\end{exercise}

\begin{exercise}
Let $\mu$ and $\nu$ be positive $\sigma$-finite measures with $\nu \ll \mu$. Show
that $\abs{\nu}(A) = \int_A \abs{\frac{d\nu}{d\mu}}\, d\mu$.
\end{exercise}

\begin{exercise}[Chain rule]
Let $\lambda$ be Lebesgue measure,
$\mu(A) = \int_A e^{-x} \ind_{[0,\infty)}(x)\, d\lambda(x)$, and
$\nu(A) = \int_A x e^{-x} \ind_{[0,\infty)}(x)\, d\lambda(x)$.
Compute $\frac{d\nu}{d\mu}$ and $\frac{d\mu}{d\lambda}$. Verify the chain rule.
\end{exercise}

\begin{exercise}
Let $(\Omega, \mathcal{F}, \PP)$ be a probability space and $X$ a random variable
with $\E[\abs{X}] < \infty$. Let $\mathcal{G} = \sigma(Y)$ where $Y$ is a discrete
random variable taking values $y_1, y_2, \ldots$ Show that:
\[
    \E[X | \mathcal{G}](\omega)
    = \sum_k \frac{\E[X \ind_{\{Y = y_k\}}]}{\PP(Y = y_k)}\, \ind_{\{Y = y_k\}}(\omega).
\]
\end{exercise}

\begin{exercise}
Prove the uniqueness of the Lebesgue decomposition: if $\nu = \nu_a + \nu_s$ and
$\nu = \nu_a' + \nu_s'$ with $\nu_a, \nu_a' \ll \mu$ and
$\nu_s, \nu_s' \perp \mu$, then $\nu_a = \nu_a'$ and $\nu_s = \nu_s'$.
\end{exercise}
